% ============================================================================
% Fichier  : partie5/P05_C23_continuite.tex
% Titre    : Continuité et Reprise
% Auteur   : MINKA MI NGUIDJOÏ Thierry Emmanuel
% Version  : 1.0 -- novembre 2026
% ============================================================================
% Base : brouillon « Livre VIII » (sécurité des opérations). Sources :
% ISO 22301:2019 ; ISO/IEC 27031:2011 ; NIST SP 800-34 Rev. 1 (2010).
% Chiffres de la reprise de LiZbiZ : cas fictif, cohérents entre eux.
% ============================================================================

\chapter{Continuité et Reprise}
\label{chap:continuite}
\chaptermark{Continuité et reprise}

\epigraph{%
    « Une sauvegarde qui n'a jamais été restaurée\\
    est une hypothèse, pas une sauvegarde. »%
}{Principe du cours}

\niveauChapitre{Intermédiaire}{Chapitres~\ref{chap:risques},
\ref{chap:architecture} et~\ref{chap:operations}.}

\begin{boxobjectifs}
\begin{itemize}
    \item Conduire une analyse d'impact sur l'activité et fixer des
          objectifs de reprise.
    \item Concevoir une stratégie de sauvegarde résistante aux
          rançongiciels.
    \item Distinguer plan de continuité et plan de reprise informatique.
    \item Choisir une stratégie de secours.
    \item Tester un plan et mesurer les écarts.
\end{itemize}
\end{boxobjectifs}

\bigskip

% ============================================================================
\section*{Question directrice : pourquoi a-t-il fallu onze jours pour facturer à nouveau~?}
% ============================================================================

LiZbiZ a repris sa facturation le 16~mars, onze jours après l'attaque. Les
sauvegardes en ligne ayant été détruites, il a fallu repartir d'une copie
de fin janvier retrouvée sur un disque oublié, puis ressaisir six semaines
de factures à partir des bordereaux papier. Personne n'avait fixé combien
de jours la facturation pouvait s'arrêter, ni combien de données on
acceptait de perdre.

\medskip
\textbf{Plan du chapitre.}
\begin{enumerate}
    \item L'analyse d'impact sur l'activité (section~\ref{sec:c23-bia}).
    \item Les sauvegardes (section~\ref{sec:c23-sauvegardes}).
    \item Continuité et reprise (section~\ref{sec:c23-plans}).
    \item Les stratégies de secours (section~\ref{sec:c23-secours}).
    \item Tester (section~\ref{sec:c23-tester}).
\end{enumerate}

\separateur

% ============================================================================
\section{L'Analyse d'Impact sur l'Activité}
\label{sec:c23-bia}
% ============================================================================

\begin{boxdef}
\termecle{Analyse d'impact sur l'activité} (BIA)~: étude qui détermine,
pour chaque activité, l'impact de son interruption dans le temps, et en
déduit la \termecle{durée maximale d'interruption admissible} (DMIA, ou
\emph{RTO}) et la \termecle{perte de données maximale admissible}
(PDMA, ou \emph{RPO})~\cite{iso22301_2019}.
\end{boxdef}

\begin{tableaucours}{L{3.4cm}L{2.0cm}L{2.0cm}L{4.8cm}}{%
    \tableauHeader{Activité} &
    \tableauHeader{RTO} &
    \tableauHeader{RPO} &
    \tableauHeader{Justification}
}
Facturation & 48~h & 4~h & Trésorerie~; pénalités des clients au-delà de
deux jours \\
Suivi des expéditions & 24~h & 1~h & Livraisons bloquées, clients sans
information \\
Paiement des fournisseurs & 5~jours & 24~h & Échéances hebdomadaires \\
Messagerie & 24~h & 24~h & Communication avec les clients \\
\end{tableaucours}
\vspace{-0.8em}
\begin{center}{\small\itshape Tableau~23.1. Objectifs de reprise de LiZbiZ
(fixés en avril, après l'attaque)}\end{center}

\begin{boxscenario}{Opération GAROUA. Les écarts de mars}
\textbf{Constatations.} Durée d'arrêt de la facturation~: du 5~mars
21~h~40 au 16~mars, soit plus de dix jours, pour un objectif fixé
après coup à 48~heures. Données perdues~: depuis la copie de fin janvier,
soit environ six semaines, pour un objectif de 4~heures.

\medskip
\textbf{Analyse} (confiance élevée). Les écarts sont de l'ordre d'un
facteur cinq sur la durée et de plusieurs centaines sur les données. Ils
mesurent l'absence de stratégie de sauvegarde et de reprise, plus encore
que l'attaque elle-même.
\end{boxscenario}

% ============================================================================
\section{Les Sauvegardes}
\label{sec:c23-sauvegardes}
% ============================================================================

\begin{boxdef}
\termecle{Règle 3-2-1}~: au moins trois copies des données, sur deux
supports différents, dont une hors site. Contre les rançongiciels, on y
ajoute une copie \emph{hors ligne ou immuable}, et la vérification
régulière de la restauration.
\end{boxdef}

\begin{tableaucours}{L{3.6cm}L{8.6cm}}{%
    \tableauHeader{Exigence} &
    \tableauHeader{Pourquoi}
}
Fréquence alignée sur le RPO & Un RPO de 4~h exige au moins une sauvegarde
toutes les 4~h (ou une réplication) \\
Copie immuable ou hors ligne & L'attaquant administrateur ne peut pas
l'effacer \\
Comptes et réseau séparés & La sauvegarde ne dépend pas du domaine
compromis (Chapitre~\ref{chap:architecture}) \\
Chiffrement et gestion des clés & La copie hors site ne fuit pas~; la clé
n'est pas perdue avec le site \\
Tests de restauration & Seule preuve que la sauvegarde sert \\
\end{tableaucours}
\vspace{-0.8em}
\begin{center}{\small\itshape Tableau~23.2. Les exigences d'une stratégie de sauvegarde}\end{center}

\begin{boxalert}
\textbf{Violation typique.} Prendre la réplication pour une sauvegarde.

\smallskip
\textbf{Règle~:} la réplication copie aussitôt les changements, y compris
le chiffrement par un rançongiciel ou une suppression par erreur. Elle
protège contre la panne d'un site, pas contre la corruption des données.
Il faut des versions conservées dans le temps.
\end{boxalert}

% ============================================================================
\section{Continuité et Reprise}
\label{sec:c23-plans}
% ============================================================================

\begin{tableaucours}{L{3.4cm}L{4.4cm}L{4.4cm}}{%
    \tableauHeader{Plan} &
    \tableauHeader{Objet} &
    \tableauHeader{Contenu}
}
Plan de continuité d'activité (PCA) & Continuer les activités essentielles
pendant la crise & Modes dégradés, organisation, communication \\
Plan de reprise informatique (PRA) & Rétablir les systèmes & Ordre de
reprise, procédures techniques, responsables \\
Plan de gestion de crise & Décider et communiquer & Cellule de crise,
contacts, messages \\
\end{tableaucours}
\vspace{-0.8em}
\begin{center}{\small\itshape Tableau~23.3. Les plans de continuité}\end{center}

Le NIST SP~800-34 décrit la planification de la contingence des systèmes
d'information~\cite{nist80034}~; l'ISO/IEC~27031 traite de la préparation
des technologies à la continuité~\cite{iso27031_2011}. Un mode dégradé
prévu (facturer sur un modèle de tableur avec une numérotation réservée,
par exemple) permet de tenir pendant la reprise.

% ============================================================================
\section{Les Stratégies de Secours}
\label{sec:c23-secours}
% ============================================================================

\begin{tableaucours}{L{2.8cm}L{4.6cm}L{2.4cm}L{2.4cm}}{%
    \tableauHeader{Stratégie} &
    \tableauHeader{Principe} &
    \tableauHeader{RTO typique} &
    \tableauHeader{Coût}
}
Site froid & Locaux et alimentation, matériel à installer & Semaines &
Faible \\
Site tiède & Matériel installé, données à restaurer & Jours & Moyen \\
Site chaud & Systèmes prêts, données répliquées & Heures & Élevé \\
Nuage & Ressources mises en route à la demande & Heures à jours & Selon
l'usage \\
\end{tableaucours}
\vspace{-0.8em}
\begin{center}{\small\itshape Tableau~23.4. Stratégies de secours}\end{center}

LiZbiZ dispose de quatre centres~: avec une architecture adaptée, chacun
peut servir de secours à un autre, à condition qu'un rançongiciel ne puisse
pas atteindre les quatre à la fois.

% ============================================================================
\section{Tester}
\label{sec:c23-tester}
% ============================================================================

\begin{tableaucours}{L{3.4cm}L{8.8cm}}{%
    \tableauHeader{Type de test} &
    \tableauHeader{Contenu}
}
Revue documentaire & Le plan est-il à jour (contacts, systèmes)~? \\
Exercice sur table & Les responsables déroulent un scénario autour d'une
table \\
Test technique partiel & Restauration d'un serveur, chronométrée \\
Simulation complète & Bascule réelle sur le site de secours \\
\end{tableaucours}
\vspace{-0.8em}
\begin{center}{\small\itshape Tableau~23.5. Les tests du plan}\end{center}

\begin{boxobjet}{M-23 · Test trimestriel de restauration}
\textbf{Objectif~:} prouver que le RTO et le RPO de la facturation sont
tenus. \textbf{Mise en œuvre~:} chaque trimestre, restauration de
l'application de facturation sur un environnement isolé, à partir de la
copie immuable. \textbf{Mesures~:} durée de restauration (comparée au
RTO), âge des données restaurées (comparé au RPO), anomalies.
\textbf{Preuve~:} procès-verbal signé, présenté au comité de sécurité.
\end{boxobjet}

\begin{boxmetier}{Administrateur systèmes et réseaux}
\textbf{Chronométrer.} L'administrateur mesure la durée réelle de chaque
restauration et la compare au RTO~; un écart déclenche une révision de la
stratégie, pas une révision de l'objectif.
\end{boxmetier}

\begin{boxmetier}{RSSI}
\textbf{Faire fixer les objectifs par les métiers.} Le RSSI fait fixer RTO
et RPO par les responsables d'activité, au regard de l'impact, et présente
à la direction le coût de chaque niveau d'objectif.
\end{boxmetier}

\separateur

\begin{boxresume}
\textbf{Ce qu'il faut retenir de ce chapitre~:}
\begin{itemize}
    \item \textbf{L'analyse d'impact} fixe, par activité, la durée
          d'interruption (RTO) et la perte de données (RPO) admissibles.
    \item \textbf{La règle 3-2-1}, complétée d'une copie immuable ou hors
          ligne, résiste au rançongiciel~; la réplication n'est pas une
          sauvegarde.
    \item \textbf{PCA, PRA et plan de crise} répondent à trois besoins~:
          continuer, rétablir, décider.
    \item \textbf{Les stratégies de secours} vont du site froid au site
          chaud, avec des coûts croissants.
    \item \textbf{Un plan se teste}~; une restauration se chronomètre.
\end{itemize}
\end{boxresume}

\begin{boxinfo}
\textbf{Points de contrôle de l'élève.} Je sais~: définir RTO et RPO et
les fixer à partir de l'impact~; appliquer la règle 3-2-1~; distinguer
réplication et sauvegarde~; distinguer PCA, PRA et plan de crise~; choisir
une stratégie de secours~; organiser un test de restauration.
\end{boxinfo}

\bigskip

\begin{boxexo}[title={\ding{45}\quad Exercice~23.1. RTO et RPO \niveaubase}]
Expliquez la différence entre RTO et RPO avec l'exemple de la facturation
de LiZbiZ.
\end{boxexo}

\begin{boxexo}[title={\ding{45}\quad Exercice~23.2. La fréquence \niveaubase}]
Avec une sauvegarde chaque nuit à 2~h, quelle est la perte de données
maximale en cas d'incident à 23~h~59~? Le RPO de la facturation est-il
tenu~?
\end{boxexo}

\begin{boxexo}[title={\ding{45}\quad Exercice~23.3. 3-2-1 \niveaubase}]
Une entreprise garde ses données sur son serveur et une copie sur un
disque branché en permanence au même serveur. Combien de règles du 3-2-1
respecte-t-elle~?
\end{boxexo}

\begin{boxexo}[title={\ding{45}\quad Exercice~23.4. Les écarts \niveaumoyen}]
Calculez, en heures, la durée d'arrêt de la facturation de LiZbiZ en mars
(en retenant une reprise le 16~mars à 8~h), et le rapport avec le RTO.
\end{boxexo}

\begin{boxexo}[title={\ding{45}\quad Exercice~23.5. Réplication \niveaumoyen}]
La BANQUE SANAGA réplique sa base en temps réel vers un second site.
Est-elle protégée contre un rançongiciel~? Que faut-il ajouter~?
\end{boxexo}

\begin{boxexo}[title={\ding{45}\quad Exercice~23.6. Mode dégradé \niveaumoyen}]
Proposez un mode dégradé pour le suivi des expéditions de LiZbiZ pendant
trois jours sans application.
\end{boxexo}

\begin{boxexo}[title={\ding{45}\quad Exercice~23.7. L'ordre de reprise \niveaumoyen}]
Dans quel ordre rétablir~: l'annuaire, la base de données, l'application
de facturation, la messagerie, le portail de suivi~? Justifiez par les
dépendances et les RTO.
\end{boxexo}

\begin{boxexo}[title={\ding{45}\quad Exercice~23.8. Le procès-verbal \niveaumoyen}]
Rédigez le procès-verbal d'un test de restauration de la facturation,
avec les mesures attendues.
\end{boxexo}

\begin{boxexo}[title={\ding{45}\quad Exercice~23.9. Stratégie \niveauexpert}]
Proposez à LiZbiZ une stratégie de sauvegarde et de secours qui tienne
les objectifs du Tableau~23.1, en utilisant ses quatre centres.
\end{boxexo}

\begin{boxexo}[title={\ding{45}\quad Exercice~23.10. L'exercice \niveauexpert}]
Concevez un exercice de reprise pour l'ANPP~: scénario, participants,
déroulé, mesures, critères de réussite.
\end{boxexo}

\bigskip

\begin{boxinfo}
\textbf{Pour aller plus loin.}
\begin{itemize}
    \item ISO~22301~\cite{iso22301_2019}~; ISO/IEC~27031~\cite{iso27031_2011}.
    \item NIST SP~800-34~\cite{nist80034}.
\end{itemize}
\end{boxinfo}

\begin{boxinfo}
\textbf{Transition.} Reprendre l'activité ne dispense pas d'établir les
faits. Le Chapitre~\ref{chap:investigation} présente les premiers gestes
de l'investigation numérique.
\end{boxinfo}
